\cotefile{63}
\scanpage{49}
\begin{flushright}
atkin \mspage{\unsure{9|8|6}}
\end{flushright}

However a significant difference arose here as to the proof of these M- and R-results. If we define for any power series
\[
\handwritten{A(x)=\sum_{n\geqslant N}a(n)x^{n}}
\]
the operator $U=U_{q}$ by
\[
\handwritten{U A(x)=\sum_{qn\geqslant N}a(qn)x^{n},}
\]
then we have
\[
\handwritten{\sum_{n=1}^{\infty}c(q^{\alpha}n)x^{n}=U^{\alpha}j(x).}
\]

As a specific example we now take the simplest case \uline{$q$}$=19$. We can show by modular theory* that there exist functions $g_{n}(x)$ for $n\handwritten{\geqslant}1$
\par\noindent\rule{\linewidth}{0.3pt}\par
\noindent *Of course the same theory could have been applied to the case $q=13$.The main difficulty is that this comparatively deep theory gives for our present purposes results which are \uline{weaker} than those obtained from an elementary use of the modular equation for $q=13$.
\par\noindent\rule{\linewidth}{0.3pt}\par
\noindent with integral coefficients\inserted{ in their power series expansions} and leading terms $x^{n}$,such that
\[
\resizebox{.86\linewidth}{!}{$\displaystyle \handwritten{\begin{aligned}[b]
 U_{19}\,j(x)&\equiv a_{1}\,g_{1}(x)+19\,a_{2}\,g_{2}(x)+19\,a_{3}\,g_{3}(x)\quad(\mathrm{mod}\;19^{2}),\\
 U_{19}\,g_{1}(x)&\equiv c_{11}\,g_{1}(x)+19\,c_{12}\,g_{2}(x)+19\,c_{13}\,g_{3}(x)\quad(\mathrm{mod}\;19^{2}),\\
 U_{19}\,g_{2}(x)&\equiv c_{21}\,g_{1}(x)+c_{22}\,g_{2}(x)+c_{23}\,g_{3}(x)+\\
 &\hspace{13em}+19\,c_{24}\,g_{4}(x)\quad(\mathrm{mod}\;19^{2}),\\
 \text{\deleted{\unsure{etc.,|etc,}}}\quad\vdots&
\end{aligned}}$}
\tag*{\handwritten{(4.1)}}
\]
\reversemarginpar\margin{\handwritten{\tikz[x=1em,y=1em,baseline=-.15em]{\draw (.57,.8)--(.07,-.7); \draw (.44,.51) .. controls (.07,.69) and (.04,.31) .. (.24,.24) .. controls (.02,.18) and (.03,.01) .. (.33,.03);}}}
where the \inserted{\handwritten{$a_{i}$ and}} $\handwritten{c_{ij}}$ are integral constants. Now it turns out \uline{by} \uline{actual} \uline{computation} that some of these constants contain powers of $19$ which were not deducible \uline{a} \uline{prio\mbox{\makebox[0pt][l]{\handwritten{\tikz[x=1em,y=1ex,baseline=0pt]{\draw (.02,-.28)--(.54,-.28)--(.54,1.54)--(1.01,1.6);}}}ir}}; moreover these full powers of $19$ which occur in actual fact are needed to prove the M- and R\normalmarginpar\margin{\handwritten{\tikz[x=1em,y=1em,baseline=-.1em]{\draw (.5,.72)--(.08,-.35) (.27,-.10)--(.38,-.35);} $-$\\(hyphen)}}\deleted{\unsure{j|;}}\inserted{\handwritten{-}}results. Thus in these cases(\uline{$q$}$=17,19,\text{and }23$)we have at present the situation that the M- and R-results can only be established for a limited number of
% archive sheet number: 48
\scanpage{50}
\begin{flushright}
at\illegible{1} \mspage{10}
\end{flushright}

\noindent fixed values of $\alpha$,and that even these limited results depend on computation beyond any reasonable prospect of hand calculation(except for $\unsure{\alpha|a}=1$). However in view of the complete test of the M-results the prggrams can be regarded as self-checking,since the M\margin{\handwritten{\tikz[x=1em,y=1em,baseline=-.1em]{\draw (.5,.72)--(.08,-.35) (.27,-.10)--(.38,-.35);} $-$\\(hyphen)}}\deleted{\unsure{j|;|,}}\inserted{\handwritten{-}}results follow theoretically from the independently computed values of $\handwritten{c_{ij}}$ in \handwritten{(4.1)} .

Since the subgroups $\handwritten{\Gamma_{0}}(13)$, $\handwritten{\Gamma_{0}}(17)$, $\handwritten{\Gamma_{0}}(19)$,and $\handwritten{\Gamma_{0}}(23)$ of the modular group have genera $0,1,1,\text{and }2$ respectively,it is remarkable that all these primes should exhibit the same behaviour with regard to our multiplicative congrue\makebox[0pt][l]{\handwritten{\tikz[x=1em,y=1ex,baseline=0pt]{\draw (.02,-.28)--(.54,-.28)--(.54,1.54)--(1.01,1.6);}}}cne p\makebox[0pt][l]{\handwritten{\tikz[x=1em,y=1ex,baseline=0pt]{\draw (.02,-.28)--(.54,-.28)--(.54,1.54)--(1.01,1.6);}}}orperties. It thus seems possible that some entirely different method or theory may exist which would give an uniform proof of all these cases.

\deleted{4}\handwritten{5}. \uline{The} \uline{cases} $q\handwritten{\geqslant}29$. As the reader may have anticipated from the non-appearance of $29$ in \handwritten{\S\,4},the direct analogy to \handwritten{(3.9)} and \handwritten{(3.10)} breaks down for $q=29$. Considering first the R-results modulo $29$ only,we can prove by the combination of modular theory and computing described in \handwritten{\S}\deleted{\unsure{3|5}}\handwritten{4} that
\[
\handwritten{\begin{aligned}
U_{29}\,j(x)&\equiv a_{1}\,g_{1}(x)+a_{2}\,g_{2}(x)\quad(\mathrm{mod}\;29),\\
U_{29}\,g_{1}(x)&\equiv 26\,g_{1}(x)+14\,g_{2}(x)\quad(\mathrm{mod}\;29),\\
U_{29}\,g_{2}(x)&\equiv 3\,g_{1}(x)+21\,g_{2}(x)\quad(\mathrm{mod}\;29).
\end{aligned}}
\]
Thus
\[
\resizebox{\linewidth}{!}{$\displaystyle \handwritten{U_{29}\{g_{1}(x)+\mu\,g_{2}(x)\}\equiv(3\mu-3)\,g_{1}(x)+(14-8\mu)\,g_{2}(x)\quad(\mathrm{mod}\;29),}$}
\]
and if we attempt to make the RHS a multiple of $\handwritten{\{}g_{\handwritten{1}}(x)\ \deleted{\unsure{\text{and }g_{2}|\text{and }\mu_{2}}}\handwritten{{}+\mu\,g_{2}(x)\}}$ we are led to the congruence
\[
\handwritten{\mu^{2}-8\mu+5\equiv 0\;(\mathrm{mod}\;29)}
\]
which has in fact no roots. However if we now consider our various coefficients \inserted{\handwritten{not}} as integers modulo $29$,but as elements of the Galois field $\mathrm{GF}(29^{2})$ we may assert(using ``$=$'' rather than ``$\handwritten{\equiv}$'' in a field of characteristic $29$):
\normalmarginpar
% archive sheet number: 49
\scanpage{51}
\begin{flushright}
atkin \mspage{11}
\end{flushright}

\[
c(29n)=a\,\handwritten{\lambda}(n)+\mathord{\rlap{$\handwritten{\bar{\phantom{a}}}$}a}\,\handwritten{\bar{\lambda}}(n),
\]
\uline{where} $a,\ \handwritten{\lambda}(n)$,\uline{and} \uline{their} \uline{conjugates} $\mathord{\rlap{$\handwritten{\bar{\phantom{a}}}$}a},\ \handwritten{\bar{\lambda}}(n)$ \uline{lie} \uline{in} $\rlap{\smash{\handwritten{\tikz[x=1em,y=1ex,baseline=0pt]{\draw (.09,-.42) .. controls (.48,-.71) and (.92,-.63) .. (1.1,-.28); \draw (.16,-.94)--(.20,-.67) (.03,1.17)--(.20,1.47); \node[inner sep=0pt] at (.87,1.95){\tiny\unsure{?}}; \draw (1.17,2.20)--(1.30,2.58);}}}}\unsure{\mathbb{GF}|\mathrm{GF}}(29^{2})$,\margin{\handwritten{\tikz[baseline=(f.base)]{\node[inner sep=1.5pt](f){$GF$};\draw (f.center) ellipse [x radius=1.12em,y radius=.78em];}}}
\[
\text{\uline{and} }\handwritten{\lambda}(29n)=\handwritten{\lambda}(29)\handwritten{\lambda}(n)\ ,\
\handwritten{\bar{\lambda}}(29n)=\handwritten{\bar{\lambda}}(29)\handwritten{\bar{\lambda}}(n).
\tag*{(5.2)}
\]

Of course $c(29n)$ itself lies in $\mathrm{GF}(29)$. We may infer from (5.2) that \handwritten{,} for all $n\handwritten{\geqslant}1$\handwritten{,}
\[
c(29^{29^{2}}n)\handwritten{\equiv}c(29n)\qquad(\mathrm{mod}\;29).
\]
We have also an M-result for $\handwritten{\lambda}(n)$ of the usual kind:
\[
\handwritten{\lambda}(np)-\handwritten{\lambda}(n)\handwritten{\lambda}(p)+p^{-1}\handwritten{\lambda}(n/p)=0\quad\text{\uline{if} }p\mathrel{\rlap{$\handwritten{/}$}=}29.
\tag*{(5.3)}
\]

It is now quite \unsure{éasy|easy} to foresee the natural general conjecture for powers of $\unsure{29'|29}$ in the form:

\uline{For} \uline{all} $\alpha\handwritten{\geqslant}1$ \uline{we} \uline{have}
\[
c(29^{\alpha}n)=a\,\handwritten{\lambda}(n)+\mathord{\rlap{$\handwritten{\bar{\phantom{a}}}$}a}\,\handwritten{\bar{\lambda}}(n),
\]
\uline{where} $a,\handwritten{\lambda}(n)$,\uline{and their conjugates} $\mathord{\rlap{$\handwritten{\bar{\phantom{a}}}$}a},\handwritten{\bar{\lambda}}(n)$,\uline{lie in an} \uline{extension} \uline{of} \uline{the} \uline{residue} \uline{class} \uline{ring} \uline{modulo} $29^{\alpha}$ \uline{defined} \uline{by} \uline{an} \uline{irreducible} \uline{quadratic} \uline{equation},\uline{and} \uline{for} $n\handwritten{\geqslant}1$ \uline{we} \uline{have}
\[
\resizebox{.86\linewidth}{!}{$\displaystyle \begin{aligned}[b]
&\handwritten{\lambda}(29n)=\handwritten{\lambda}(29)\handwritten{\lambda}(n)\\
\text{\uline{and} }\;&\handwritten{\lambda}(np)-\handwritten{\lambda}(n)\handwritten{\lambda}(p)+p^{-1}\handwritten{\lambda}(n/p)=0\quad\text{if }p\mathrel{\rlap{$\handwritten{/}$}=}29.
\end{aligned}$}
\tag*{(5.4)}
\]

This is subject to the same reservations as the cases $q=17,19,\text{and }23$; it is only proved for $\alpha\handwritten{\leqslant}5$,and then only by a combination of modular theory and computation.

The case $q=31$ is precisely analogous to $q=29$,while for $q=37$ we find that \uline{three} conjugates are required over an extension of the ring modulo $37^{\alpha}$ defined by an irreducible \uline{cubic} equation.At this point it seemed unprofitable to compute any further cases in detail, since no essentially new features arose after $29$. However brief pre\inserted{\handwritten{l}}iminary investigations for all primes \uline{$q$} with $41\handwritten{\leqslant}q\handwritten{\leqslant}67$ suggest that in general $c(q^{\alpha}n)$ will be a linear combination of $\handwritten{[}q/12\handwritten{]}$ \inserted{\handwritten{multiplicative}} coefficients lying in various extensions of the ring modulo $q^{\alpha}$. In addition the analogy with Hecke theory for modular forms of negative dimension
% archive sheet number: 50
\scanpage{52}
\begin{flushright}
atkin \mspage{12}
\end{flushright}

\noindent which we now discuss \unsure{bríefly|briefly} in \handwritten{\S}\,6 below indicates that (5.4) is the right kind of property to generalise for all \uline{$q$}.

6. \uline{The} \uline{analogy} \uline{with} \uline{Hecke} \uline{theory}. We now return to the exact multiplicative properties such as (3.5). These have been extensively generalised by Hecke,Petersson,and others; for a useful summary we refer to RANKIN(196\hspace{.6em}). As a particular example of the theory we take the case of
\[
\handwritten{\sum_{n=2}^{\infty}}t_{2}(n)x^{n}=x^{2}f^{48}(x),
\]
for which Hecke showed that
\[
t_{2}(n)=a\,\handwritten{\lambda}(n)+\mathord{\rlap{$\handwritten{\bar{\phantom{a}}}$}a}\,\handwritten{\bar{\lambda}}(n),
\]
where $a,\handwritten{\lambda}(n)$,and their conjugates $\mathord{\rlap{$\handwritten{\bar{\phantom{a}}}$}a},\handwritten{\bar{\lambda}}(n)$ lie in the quadratic field $\deleted{\unsure{q|Q}}$\reversemarginpar\margin{\handwritten{\tikz[x=1em,y=1em,baseline=-.1em]{\draw (.56,.7)--(.03,-.53) (.78,.39)--(.20,-.74); \draw (.33,.39) .. controls (.13,.46) and (.14,.16) .. (.27,.04) .. controls (.4,-.11) and (.26,-.26) .. (.25,-.3);}}} $Q(\handwritten{\surd}\,144169)$,and for $n\handwritten{\geqslant}1$ and prime $p$ we have
\[
\handwritten{\lambda}(np)-\handwritten{\lambda}(n)\handwritten{\lambda}(p)+p^{23}\handwritten{\lambda}(n/p)=0.
\tag*{(6.1)}
\]

It seems clear that the multiplicative congruence properties we have described in terms of extensions of residue class rings are closely anal\deleted{g}\inserted{\handwritten{\unsure{a|u|o}}}gous to the exact multiplicative properties of Hecke theory in terms of algebraic number fields. We have not described in this account the underlying common feature of the Hecke operators,but it seems reasonable to conjecture that some kind of ``$\handwritten{\mathfrak{p}}$-adic'' Hecke theory\normalmarginpar\margin{\handwritten{$\left(\begin{gathered}\text{German}\\[-.4ex]p\end{gathered}\right)$}} exists for modular forms of positive and(as for $c(n)$) zero dimension, which is a counterpart of the exact theory for forms of negative dimension. The present difficulties in the way of precise formulation and proof see\inserted{\handwritten{m}} formidable in so far as the \uline{congruences} we have considered in this paper arise essentially at just one point($\inserted{\handwritten{\tau=}}i\handwritten{\infty}$) of the complex upper half-plane,whereas the \uline{identities} of Hecke theory exist at all points of the upper half-plane,taking the form of multiplicative properties for Fourier coefficients at $\inserted{\handwritten{\tau=}}i\handwritten{\infty}$.

7. \uline{Non-Ramanujan} \uline{congruences} \uline{for} \uline{primes} $q\handwritten{\leqslant}11$. Once the above properties for primes $q\handwritten{\geqslant}13$ had been found it was clear that something similar must
% archive sheet number: 51
\scanpage{53}
\begin{flushright}
at\illegible{1} \mspage{13}
\end{flushright}

\noindent exist for the smaller primes as well. As an example with $q=5$ we find:

\uline{Let}*$t(n)=c(5^{\alpha}n)/c(5^{\alpha})$. \uline{Then}
\par\noindent\rule{\linewidth}{0.3pt}\par
\noindent * Since\reversemarginpar\margin{\handwritten{\unsure{online|on line}}} $\handwritten{\mathord{\tikz[x=1ex,y=1ex,baseline=0pt]{\draw (.82,1.85) .. controls (.42,1.55) and (.34,-.54) .. (0,-.35);}}}\raisebox{.37ex}{$5^{\alpha+1}$}\handwritten{\mathord{\tikz[x=1ex,y=1ex,baseline=0pt]{\draw (0,1.62) .. controls (.04,1.02) and (-.13,.29) .. (.35,.1);}}}\handwritten{\mid}c(5^{\alpha}n)$ and $5^{\alpha+1}\handwritten{\parallel}c(5^{\alpha})$ this definition is legitimate in any residue class ring modulo $5^{\beta}$.
\par\noindent\rule{\linewidth}{0.3pt}\par
\[
\resizebox{.86\linewidth}{!}{$\displaystyle \begin{aligned}[b]
&t(5n)\handwritten{\equiv}t(5)t(n)\qquad(\mathrm{mod}\;5^{3\alpha}),\\
\text{\uline{and} }\;&t(np)-t(n)t(p)+p^{-1}t(n/p)\handwritten{\equiv}0\quad(\mathrm{mod}\;5^{3\alpha})\quad\text{\uline{if} }p\mathrel{\rlap{$\handwritten{/}$}=}5\ \text{\uline{is} \uline{prime}.}
\end{aligned}$}
\tag*{(7.1)}
\]

For $q=2,3,5,\text{and }7$ these results can be fully proved as for $q=13$; for $q=11$ we have again the same reservations as for $q=19$. In a s\makebox[0pt][l]{\handwritten{\tikz[x=1em,y=1ex,baseline=0pt]{\draw (.02,-.28)--(.54,-.28)--(.54,1.54)--(1.01,1.6);}}}nese we may feel that the Ramanujan-type properties (2.5) which do not generalise to higher primes have previously concealed the existence of the multiplicative properties. Thus(to suppose what is in fact false) had $c(2^{\alpha}\inserted{\handwritten{\text{\uline{$n$}}}})$ been prime to $2$,its multiplicative properties might have been sought and found by hand; as it is\handwritten{,}the first non-Ramanujan prime is $13$ which is beyond the scope of hand computation. On the other hand had $p(5n+4)$ been prime to $5$ the whole possibility of considering congruences for modular forms m\inserted{\handwritten{i}}ght never have been noticed.

We may finally try to summarise the distinction between Ramanujan-type and multiplicative properties in this way. The Ramanujan-type property asserts that
\[
\handwritten{\sum_{n=1}^{\infty}}c(5^{\alpha}n)\inserted{\handwritten{x^{n}}}=5^{\alpha+1}P(x),
\]
where \uline{$P(x$}$)$ is a power series with integral coefficients. The multiplicative property now asserts that \uline{$P(x$}$)$ itself is an interesting series.

8. \uline{Notes} \uline{on} \uline{computation}. The main requirement for most of this work was the ability to compute the Fo\makebox[0pt][l]{\handwritten{\tikz[x=1em,y=1ex,baseline=0pt]{\draw (.02,-.28)--(.54,-.28)--(.54,1.54)--(1.01,1.6);}}}ruier expansions of various modular forms and functions modulo \uline{$m$}. If we have available in store various power series $g_{i}(x)$ up to the term in $x^{N}$,then it is clear that addition and subtraction of constant multiples of these series,testing for ``accidents'',etc.,takes a time approximately proportional to $N$; we shall abbreviate this to ``takes $\unsure{O|0}(N)$''. The test for M-results takes
\normalmarginpar
% archive sheet number: 52
\scanpage{54}
\begin{flushright}
\illegible{1}\unsure{t}kin \mspage{14}
\end{flushright}

\[
\handwritten{O\left(\sum_{p<N}N/p\right)=O\left(N\unsure{-\log|\log}-\log N\right)}
\]
$=\unsure{O|0}(N)$ in practice for computers in this millenium.

Next if we have a given series $g(x)$,and multiply(divide) it by Euler's series*

\noindent\rule{\linewidth}{0.3pt}\par
\noindent *or other suitable $\handwritten{\unsure{\vartheta|\theta|\mathfrak{g}}}$-series of the form $\handwritten{\sum\lambda(n)\ldotp x^{an^{2}+\unsure{bn|b-n}}}$
\par\noindent\rule{\linewidth}{0.3pt}
{\small
\[
\resizebox{\linewidth}{!}{$\displaystyle f(x)=\handwritten{\prod_{r=1}^{\infty}}(1-x^{r})=1-x-x^{2}+x^{5}+x^{7}-x^{12}\ldots=1+\handwritten{\sum_{n=1}^{\infty}}(-1)^{n}\handwritten{\left\{x^{\frac12 n(3n-1)}+x^{\frac12 n(3n+1)}\right\},}$}
\]
}
this takes $O(N^{3/2})$. The subroutines for this were expertly coded by Dr. Paul Bryant \inserted{from my Fortran version}. Specifically the time required for $1000N$ terms was about $\handwritten{\frac14}N^{3/2}\deleted{\unsure{/4|/-4}}$\reversemarginpar\margin{\handwritten{\tikz[x=1ex,y=1ex,baseline=0pt]{\draw (0,-1.1)--(1.8,2.5) (.95,1.23) .. controls (1.78,2.7) and (.25,1.7) .. (.55,.6) .. controls (1.25,-.1) and (1.38,.36) .. (.91,.76);}}} seconds,so that(for example)to compute \uline{$p(n$}$)$ modulo \uline{$m$} for $n\handwritten{\leqslant}40000$ takes about a minute. The reduction modulo $m$ in the operation
\[
\handwritten{a(n)=b(n)-b(n-1)-b(n-2)+b(n-5)+b(n-7)-b(n-12)\ldots}
\]
is performed only at the end(and not after each individual $+$ or $-$ ), so that to avoid overflow there is a restriction that $m\unsure{=|\cdot}(N/6)^{1/2}$ be less than $6.7\times10^{10}$.

Finally multiplication(division) of two given general series takes $O(N^{2})$ \inserted{and had to be avoided at all costs}. In view of this a certain \unsure{amount} of modular(rather than programming)ingenuity was required; in one case a natural definition $g^{2}(x)$ with $g(x)$ already computed was replaced by $g^{2}(x)=$ a sum of three terms each involving several $\handwritten{\vartheta}$-multiplications and divisions.

The possibility of multiplying(dividing) by Jacobi's series
\[
f^{3}(x)=1-3x+5x^{3}-\ldots=\handwritten{\sum_{n=0}^{\infty}(-1)^{n}(2n+1)x^{\frac12 n(n+1)}}
\]
was considered,but it turned out to be slightly slower than 3 successive uses of Euler's series,and also involved a heavier restriction on the size of $m$ and $N$.

In the end the limiting factor for most of these programs was store rather than speed. Although the Atlas Supervisor has a ``drum learning'' program which enables many jobs to use the concept of a one-level store of about 90K without serious loss of efficiency,
\normalmarginpar
% archive sheet number: 53
\scanpage{55}
\begin{flushright}
atkin \mspage{15}
\end{flushright}

\reversemarginpar the random access involv\margin{\handwritten{/}}\deleted{\unsure{a|e}}\inserted{\handwritten{e}}d in using Euler's series to multiply (divide) \uline{$g(x$}$)$ meant that almost all the terms of \uline{$g(x$}$)$ had to be stored in core. In practice $N=43000$ was about the limit.

The main programs for each prime \uline{$q$} were written in Fortran calling on the coded subroutines described above. The only additional feature was a Fortran function $\mathrm{IJMODK}(\mathrm{K},\unsure{0|\mathrm{O}},\mathrm{I},\mathrm{J})$ which obtains the residue of $\mathrm{I}\ldotp\mathrm{J}\ (\mathrm{mod}\ \mathrm{K})$ and overwrites its calling sequence in the\normalmarginpar\margin{\handwritten{\unsure{''}}} execution run,which Dr. Don Russell ingeniously devised for me. This allows both $\mathrm{I}$ and $\mathrm{J}$ to be as large as $6.7\times10^{10}$,whereas the direct use of Fortran would limit $\mathrm{I}\ldotp\mathrm{J}$ to this size,and also be about three times as slow. It might be appropriate here to express a hope that computer design should allow for the easy extraction of an integer \uline{\unsure{remain\`{d}er|remainder}} after division,as \reversemarginpar w\margin{\handwritten{/}}\deleted{a}\inserted{\handwritten{e}}ll as the quotient mainly required in floating point work.

Finally it may be of interest that some of these programs when run in 1964 exposed a low-probability \inserted{timing} fault in the machine's hardware. Whereas a convergent process of numerical integration(for example) may not be aff\deleted{\unsure{t|b}}\unsure{e|o}cted if a fault occurs once in a million times, arithmetic modulo \uline{$m$} depending on remainders rather than quotients is totally destroyed.
\normalmarginpar
% archive sheet number: 54
